% =====================================================================
%  CARNET D'ENTRAINEMENT — FICHE 6 — THÈME 3
%  Niveau MOYEN — Corrigé
% =====================================================================
\documentclass[11pt,a4paper]{article}
\usepackage[utf8]{inputenc}
\usepackage[T1]{fontenc}
\usepackage{lmodern}
\usepackage[french]{babel}
\usepackage{amsmath,amssymb,mathtools}
\usepackage{icomma}
\usepackage[version=4]{mhchem}
\usepackage{siunitx}
\sisetup{output-decimal-marker={,}}
\usepackage{xcolor}
\usepackage{colortbl}
\usepackage{tikz}
\usepackage[most]{tcolorbox}
\usepackage{enumitem}
\usepackage{array}
\usepackage{booktabs}
\usepackage{fancyhdr}
\usepackage[hidelinks]{hyperref}
\usetikzlibrary{arrows.meta,positioning,calc,shapes.geometric}
\usepackage[a4paper,margin=2.1cm,headheight=26pt,headsep=8pt]{geometry}

\definecolor{primary}{RGB}{146,95,160}
\definecolor{primarydark}{RGB}{92,52,104}
\definecolor{corcol}{RGB}{0,114,178}
\definecolor{astucecol}{RGB}{198,93,9}

\tcbset{boxbase/.style={enhanced,breakable,boxrule=0.9pt,arc=2.5pt,
   left=3mm,right=3mm,top=2.2mm,bottom=2.2mm,
   fonttitle=\bfseries,coltitle=white,toptitle=0.6mm,bottomtitle=0.6mm}}
\newtcolorbox[auto counter]{cor}[1][]{boxbase,colback=corcol!5,colframe=corcol,
   title={\textsc{Corrigé}~\thetcbcounter\ifx\relax#1\relax\relax\else\textmd{ —\itshape\ #1}\fi}}
\newtcolorbox{astucebox}[1][]{boxbase,colback=astucecol!8,colframe=astucecol,
   title={\textsc{Astuce}\ifx\relax#1\relax\relax\else\textmd{ : \itshape #1}\fi}}

\pagestyle{fancy}\fancyhf{}
\fancyhead[L]{\small\textcolor{primarydark}{\textbf{Fiche 6 · Thème 3} — Corrigé}}
\fancyhead[R]{\small\textcolor{primarydark}{\textsc{Moyen}}}
\fancyfoot[C]{\small\thepage}
\renewcommand{\headrulewidth}{0.8pt}
\renewcommand{\headrule}{\hbox to\headwidth{\color{primary}\leaders\hrule height \headrulewidth\hfill}}
\setlength{\parindent}{0pt}\setlength{\parskip}{3pt}

\newcommand{\rep}{\textbf{\textcolor{corcol}{Réponse : }}}

\begin{document}
\begin{center}
{\Large\bfseries\color{primarydark}Thème 3 — Corrigé du niveau Moyen}
\end{center}
\medskip

\begin{cor}[ion oxonium \ce{H3O+}]
\begin{enumerate}[label=\alph*),leftmargin=2em,itemsep=2pt]
  \item Oxygène : $3$ liaisons, $1$ doublet libre.\;
        $q_{\mathrm{f}}(\ce{O}) = 6 - 2\times 1 - 3 = 6 - 2 - 3 = \mathbf{+1}$.
  \item Chaque hydrogène : $1$ liaison, $0$ doublet libre.\;
        $q_{\mathrm{f}}(\ce{H}) = 1 - 0 - 1 = \mathbf{0}$.
  \item Somme $= (+1) + 3\times 0 = \mathbf{+1}$. \checkmark\;
        \rep Cohérent avec la charge $+1$ de l'ion \ce{H3O+} : c'est l'oxygène, qui a « gagné »
        une troisième liaison, qui porte la charge positive.
\end{enumerate}
\end{cor}

\begin{cor}[ion carbonate \ce{CO3^2-}]
\begin{enumerate}[label=\alph*),leftmargin=2em,itemsep=2pt]
  \item Carbone : $0$ doublet libre, $4$ liaisons $\Rightarrow q_{\mathrm{f}} = 4 - 0 - 4 = \mathbf{0}$.\\
        \ce{O} doublement lié : $2$ doublets libres, $2$ liaisons $\Rightarrow
        q_{\mathrm{f}} = 6 - 2\times 2 - 2 = \mathbf{0}$.\\
        \ce{O} simplement lié : $3$ doublets libres, $1$ liaison $\Rightarrow
        q_{\mathrm{f}} = 6 - 2\times 3 - 1 = \mathbf{-1}$ (pour chacun des deux).
  \item Somme $= 0 + 0 + (-1) + (-1) = \mathbf{-2}$. \checkmark\; C'est bien la charge de l'ion.
  \item \rep \textbf{Oui}, les charges négatives sont bien placées : elles se trouvent sur les
        \textbf{oxygènes}, plus \emph{électronégatifs} que le carbone. Une charge négative est
        d'autant plus stable qu'elle est portée par l'atome le plus électronégatif.
\end{enumerate}
\end{cor}

\begin{cor}[ion ammonium \ce{NH4+}]
\begin{enumerate}[label=\alph*),leftmargin=2em,itemsep=2pt]
  \item Azote : $0$ doublet libre, $4$ liaisons $\Rightarrow q_{\mathrm{f}}(\ce{N}) = 5 - 0 - 4 = \mathbf{+1}$.\;
        Chaque \ce{H} : $q_{\mathrm{f}} = 1 - 0 - 1 = \mathbf{0}$.\;
        Somme $= (+1) + 4\times 0 = \mathbf{+1}$. \checkmark
  \item \rep Normalement l'azote fait $3$ liaisons ($N_v = 5$, soit $3$ liaisons $+$ $1$ doublet
        libre). Dans \ce{NH4+}, il utilise ce doublet libre pour former une \textbf{quatrième}
        liaison (dative) avec \ce{H+}. Il passe donc à $4$ liaisons : la formule donne alors
        $q_{\mathrm{f}} = 5 - 4 = +1$. Bien qu'ayant \emph{fourni} les deux électrons de la liaison
        dative, l'azote « tétravalent » se retrouve avec un électron propre de moins qu'à l'état
        neutre : c'est lui qui porte la charge positive de l'ion.
\end{enumerate}
\end{cor}

\begin{cor}[molécule d'ozone \ce{O3}]
Structure \ce{O=O-O}.
\begin{enumerate}[label=\alph*),leftmargin=2em,itemsep=2pt]
  \item Oxygène central : la double liaison compte pour $2$ liaisons, plus la simple $=3$ liaisons,
        et $1$ doublet libre.\;
        $q_{\mathrm{f}} = 6 - 2\times 1 - 3 = 6 - 2 - 3 = \mathbf{+1}$.
  \item \ce{O} doublement lié : $2$ doublets libres, $2$ liaisons $\Rightarrow
        q_{\mathrm{f}} = 6 - 2\times 2 - 2 = \mathbf{0}$.\\
        \ce{O} simplement lié : $3$ doublets libres, $1$ liaison $\Rightarrow
        q_{\mathrm{f}} = 6 - 2\times 3 - 1 = \mathbf{-1}$.
  \item Somme $= (+1) + 0 + (-1) = \mathbf{0}$. \checkmark\;
        \rep Cohérent avec une molécule \textbf{neutre}. L'ozone porte donc une charge $+$ et une
        charge $-$ « internes » qui se compensent (édifice globalement neutre).
\end{enumerate}
\end{cor}

\begin{cor}[comparer deux structures d'un même atome]
\begin{enumerate}[label=\alph*),leftmargin=2em,itemsep=2pt]
  \item \rep La structure \textbf{(II)} est la plus représentative. Critère : à octet respecté, on
        préfère la structure aux charges formelles \textbf{minimales} (les plus proches de zéro).
        Une charge $q_{\mathrm{f}} = 0$ sur le soufre est bien plus favorable qu'une charge $+2$.
  \item \rep Le soufre est en \textbf{3\textsuperscript{e} période} : il peut faire un
        \textbf{octet étendu}. En transformant des liaisons simples en \textbf{doubles liaisons}
        avec ses voisins, il augmente son nombre de liaisons, ce qui \emph{abaisse} sa charge
        formelle de $+2$ à $0$.
\end{enumerate}
\end{cor}

\begin{astucebox}[le réflexe de la charge formelle]
$q_{\mathrm{f}} = N_v - 2\times(\text{doublets non liants}) - (\text{nb de liaisons})$, et
\textbf{une double liaison compte pour 2 liaisons}. Vérifie toujours : la somme des $q_{\mathrm{f}}$
doit redonner la charge de l'édifice.
\end{astucebox}

\end{document}
